\documentclass[A4,svgnames,9pt,aspectratio=169]{beamer}

\usepackage[utf8]{inputenc}
\usepackage[T1]{fontenc}
\usepackage[french]{babel}
\usepackage{graphicx}
\usepackage{amsmath,amssymb}
\usepackage{booktabs}
\usepackage{tabularx}
\usepackage{xcolor}
\usepackage{tikz}
\usetikzlibrary{arrows.meta,positioning}

\usetheme{inria}
\graphicspath{{figures/}}

\hypersetup{
  colorlinks=true,
  allcolors=rouge_inria,
  pdfauthor={Louis Hauseux, Raphaël Antoine, Philippe Foucher, Pierre Charbonnier, Josiane Zerubia},
  pdftitle={CrackSAM-GeoLoRA},
  pdfsubject={Guidage géométrique dans SAM 2 : résultat négatif ; gain de la dilatation},
  pdfkeywords={SAM 2, CrackSAM, LoRA, Frangi, IoU tolérante}
}

\definecolor{success}{HTML}{3C8D5A}
\definecolor{warning}{HTML}{D98200}
\definecolor{failure}{HTML}{C9191E}
\definecolor{inkblue}{HTML}{27348B}
\definecolor{softblue}{HTML}{E9F2F8}
\definecolor{softred}{HTML}{FBEAEC}
\definecolor{softgreen}{HTML}{EAF5EE}
\definecolor{softgray}{HTML}{F1F2F4}

\newcommand{\IoU}{\mathrm{IoU}}
\newcommand{\source}[1]{%
  \par\vfill{\raggedright\fontsize{5.8}{6.6}\selectfont\color{black!62}#1\par}%
}
\newcommand{\statusbadge}[2]{%
  \tikz[baseline=(n.base)]\node[rounded corners=2pt,fill=#1!13,draw=#1!55,
  inner xsep=4pt,inner ysep=1.5pt,text=#1,font=\bfseries\scriptsize] (n) {#2};%
}
\newcommand{\metriccard}[4]{%
  \begin{tikzpicture}
    \node[rounded corners=4pt,draw=#1!50,fill=#1!8,minimum width=3.05cm,
      minimum height=1.35cm,align=center,inner sep=4pt] {%
      {\scriptsize #2}\\[-1pt]
      {\Large\bfseries\color{#1} #3}\\[-1pt]
      {\fontsize{5.8}{6.5}\selectfont\color{black!62} #4}%
    };
  \end{tikzpicture}%
}
% Vignette d'un panneau, sur cinq colonnes.
\newcommand{\widepanel}[2]{%
  \begin{minipage}[t]{0.188\textwidth}
    \centering
    {\fontsize{6.1}{6.8}\selectfont\bfseries #1}\par\vspace{1pt}
    \setlength{\fboxsep}{0.7pt}%
    \fcolorbox{black!16}{white}{%
      \includegraphics[width=0.94\linewidth,height=1.86cm,keepaspectratio]{#2}}
  \end{minipage}%
}
% Vignette d'un panneau, sur quatre colonnes.
\newcommand{\compactpanel}[2]{%
  \begin{minipage}[t]{0.235\textwidth}
    \centering
    {\fontsize{6.1}{6.8}\selectfont\bfseries #1}\par\vspace{1pt}
    \setlength{\fboxsep}{0.7pt}%
    \fcolorbox{black!16}{white}{%
      \includegraphics[width=0.93\linewidth,height=1.80cm,keepaspectratio]{#2}}
  \end{minipage}%
}
\newcommand{\caserow}[6]{%
  \noindent
  \begin{minipage}[c]{0.62\linewidth}
    {\scriptsize\bfseries #1}
  \end{minipage}\hfill
  \begin{minipage}[c]{0.36\linewidth}
    \raggedleft\statusbadge{#5}{#6}
  \end{minipage}\par
  \vspace{1pt}
  \compactpanel{Image}{#2_image.jpg}\hfill
  \compactpanel{Vérité terrain}{#2_gt.jpg}\hfill
  \compactpanel{#3}{#2_ref.jpg}\hfill
  \compactpanel{#4}{#2_var.jpg}%
}
\newcommand{\threewayrow}[4]{%
  \noindent
  \begin{minipage}[c]{0.56\linewidth}
    {\scriptsize\bfseries #1}
  \end{minipage}\hfill
  \begin{minipage}[c]{0.42\linewidth}
    \raggedleft\statusbadge{warning}{#4}
  \end{minipage}\par
  \vspace{1pt}
  \widepanel{Image}{#2_image.jpg}\hfill
  \widepanel{Vérité terrain}{#2_gt.jpg}\hfill
  \widepanel{Baseline}{#2_ref.jpg}\hfill
  \widepanel{\texttt{tol3}}{#2_var.jpg}\hfill
  \widepanel{\texttt{geo\_tol3}}{#3_var.jpg}%
}
\newcommand{\synthpanel}[2]{%
  \begin{minipage}[t]{0.185\textwidth}
    \centering
    {\fontsize{6.1}{6.8}\selectfont\bfseries #1}\par\vspace{1pt}
    \setlength{\fboxsep}{0.7pt}%
    \fcolorbox{black!16}{white}{\includegraphics[width=0.86\linewidth]{#2}}
  \end{minipage}%
}
\newcommand{\colourkey}{%
  \textcolor{success}{$\blacksquare$}~vrai positif \quad
  \textcolor{failure}{$\blacksquare$}~faux positif \quad
  \textcolor{inkblue}{$\blacksquare$}~manqué}

\title[CrackSAM-GeoLoRA]{CrackSAM-GeoLoRA}
\subtitle{Le guidage géométrique n'aide pas. La dilatation, un peu.}
\date[09/08/2026]{9 août 2026}
\author[L. \textsc{Hauseux} et al.]{%
  {\small Louis \textsc{Hauseux}, Josiane \textsc{Zerubia}}\\
  {\small Raphaël \textsc{Antoine}, Philippe \textsc{Foucher}, Pierre \textsc{Charbonnier}}%
}
\institute{Inria, équipe \textsc{Ayana} \quad---\quad Cerema, équipe \textsc{Endsum}}

\begin{document}

\frame{\titlepage}

% -----------------------------------------------------------------------------
\begin{frame}{Deux résultats}
  \centering
  \vspace{0.10cm}
  \begin{tikzpicture}
    \node[rounded corners=6pt,draw=failure!50,fill=softred,text width=5.4cm,
      minimum height=3.1cm,align=left,inner sep=9pt,anchor=north west]
      at (0,0) {%
      {\large\bfseries\color{failure} Géométrie : rien}\\[6pt]
      Évidence \textbf{alignée} contre évidence \textbf{d'une autre image} :\\[4pt]
      $|\Delta| < 0{,}001$ à toutes les tolérances\\[2pt]
      le contrôle est devant \textbf{5 fois sur 6}\\[6pt]
      {\scriptsize Par image : $-0{,}0007$ en moyenne,
      603 gains, 866 pertes.}};
    \node[rounded corners=6pt,draw=success!50,fill=softgreen,text width=5.4cm,
      minimum height=3.1cm,align=left,inner sep=9pt,anchor=north west]
      at (6.4,0) {%
      {\large\bfseries\color{success} Dilatation : un peu}\\[6pt]
      Perte tolérante à \textbf{3 px} contre baseline :\\[4pt]
      \textbf{$+0{,}0035$} en IoU stricte\\[2pt]
      \textbf{$+0{,}0179$} à 1 px de tolérance\\[6pt]
      {\scriptsize Aucun paramètre ajouté, aucun coût de calcul.}};
  \end{tikzpicture}

  \vspace{0.34cm}
  \statusbadge{inkblue}{6 variantes, budget strictement égal, 1\,695 images de test}

  \source{Khánh Hà, monomodal visible. 8--9 août 2026, RTX PRO 6000 Blackwell.}
\end{frame}

% -----------------------------------------------------------------------------
\begin{frame}{Ce qu'on ajoute à la Frangi--similarité}
  \centering
  \vspace{0.04cm}
  \widepanel{Image}{b2t_gain_ct6774_image.jpg}\hfill
  \widepanel{\texttt{frangi\_sim}}{b2t_gain_ct6774_frangi.jpg}\hfill
  \widepanel{\texttt{ofs}}{b2t_gain_ct6774_ofs.jpg}\hfill
  \widepanel{\texttt{ofa}}{b2t_gain_ct6774_ofa.jpg}\hfill
  \widepanel{\texttt{phase\_sym}}{b2t_gain_ct6774_phase.jpg}

  \vspace{0.22cm}
  \small
  \begin{tabularx}{0.94\linewidth}{>{\ttfamily\footnotesize}p{4.1cm}X}
    \toprule
    \normalfont\bfseries 11 canaux, jamais multipliés &
      \normalfont\bfseries Rôle \\
    \midrule
    frangi\_sim & la carte des essais précédents : propose \\
    ofs, ofa, abs\_energy & flux orienté : récuse une frontière d'ombre \\
    even\_odd, dbic, profile & modèle explicite : ligne ou marche \\
    phase\_sym & Gabor en quadrature : polarité sombre \\
    cos2$\theta$, sin2$\theta$, scale & orientation et largeur \\
    \bottomrule
  \end{tabularx}

  \vspace{0.18cm}
  \statusbadge{failure}{filtres accordés sur 19,1 px : trop épais pour ces fissures}

  \source{\texttt{cracktree200\_6774}. Évidence calculée à $224^2$.}
\end{frame}

% -----------------------------------------------------------------------------
\begin{frame}{Où elle entre dans SAM 2}
  \centering
  \vspace{0.10cm}
  \begin{tikzpicture}[>=Latex,font=\scriptsize]
    \node[rounded corners=4pt,draw=black!30,fill=softgray,text width=2.0cm,
      align=center,inner sep=4pt] at (0,0.95) (rgb) {Image RGB\\$448^2$};
    \node[rounded corners=4pt,draw=warning!55,fill=warning!8,text width=2.35cm,
      align=center,inner sep=4pt] at (3.0,0.95) (hiera) {%
      \textbf{SAM 2 Hiera-L}\\LoRA q/v $r=4$\\453\,248 params};
    \node[rounded corners=4pt,draw=black!30,fill=softgray,text width=2.0cm,
      align=center,inner sep=4pt] at (0,-1.15) (geo) {11 canaux\\$224^2$};
    \node[rounded corners=4pt,draw=inkblue!55,fill=softblue,text width=2.35cm,
      align=center,inner sep=4pt] at (3.0,-1.15) (enc) {%
      \textbf{Encodeur géométrique}\\290\,801 params};
    \node[rounded corners=4pt,draw=inkblue!45,fill=white,text width=3.35cm,
      align=left,inner sep=4pt] at (7.1,-1.15) (inj) {%
      \textcolor{inkblue}{\textbf{3 projections init. à ZÉRO}}\\[1.5pt]
      $+$ \texttt{high\_res[0]} \hfill $32{\times}256^2$\\
      $+$ \texttt{high\_res[1]} \hfill $64{\times}128^2$\\
      $+$ \texttt{embeddings} \hfill $256{\times}64^2$};
    \node[rounded corners=4pt,draw=success!50,fill=softgreen,text width=2.35cm,
      align=center,inner sep=4pt] at (10.7,0.95) (dec) {%
      \textbf{Mask decoder}\\logits $448^2$};
    \draw[->,thick] (rgb) -- (hiera);
    \draw[->,thick] (geo) -- (enc);
    \draw[->,thick] (enc) -- (inj);
    \draw[->,thick] (hiera) -- (dec);
    \draw[->,thick] (inj.east) -| (dec.south);
  \end{tikzpicture}

  \vspace{0.30cm}
  \statusbadge{inkblue}{à l'initialisation, le modèle \emph{est} la baseline}
  \hspace{0.3cm}
  \statusbadge{success}{correction additive, révocable}
  \hspace{0.3cm}
  \statusbadge{failure}{\texttt{mask\_input} jamais utilisé}

  \source{L'invite dense de juillet coûtait $-0{,}0979$ d'IoU : cette interface
  reste un contrôle négatif.}
\end{frame}

% -----------------------------------------------------------------------------
\begin{frame}{Le protocole : six barreaux, un contrôle}
  \small
  \begin{tabularx}{\linewidth}{>{\ttfamily}p{2.75cm}p{2.3cm}>{\raggedright\arraybackslash}p{1.7cm}r>{\raggedright\arraybackslash}X}
    \toprule
    \normalfont Variante & Perte & Évidence & Params &
      \normalfont Isole \\
    \midrule
    baseline & CrackSAM & --- & 453\,248 & l'ancre \\
    cldice & $+$ clDice & --- & 453\,248 & la continuité \\
    geo & $+$ clDice & alignée & 744\,049 & la géométrie \\
    \textbf{tol3} & $+$ \textbf{tolérante 3 px} & --- & 453\,248 &
      \textbf{la dilatation} \\
    geo\_tol3 & $+$ tolérante 3 px & alignée & 744\,049 & la géométrie \\
    \textbf{geo\_tol3\_permuted} & $+$ tolérante 3 px &
      \textbf{d'une autre image} & 744\,049 & \textbf{le contrôle} \\
    \bottomrule
  \end{tabularx}

  \vspace{0.30cm}
  \centering
  \metriccard{inkblue}{Budget}{5 époques}{même graine, même ordre}
  \hspace{0.25cm}
  \metriccard{success}{Corpus}{9\,121 / 1\,695}{train / test}
  \hspace{0.25cm}
  \metriccard{warning}{Durée}{40--41 min}{par variante}

  \source{Le contrôle permuté a la même capacité et la même perte : seul
  l'appariement image~$\leftrightarrow$~évidence est détruit. Sans lui, on ne
  sépare pas la géométrie des paramètres ajoutés.}
\end{frame}

% -----------------------------------------------------------------------------
\begin{frame}{La géométrie n'aide pas}
  \begin{columns}[T,onlytextwidth]
    \begin{column}{0.42\textwidth}
      \centering
      \small
      \setlength{\tabcolsep}{3pt}%
      \begin{tabularx}{\linewidth}{Xrrr}
        \toprule
        $k$ & \ttfamily alignée & \ttfamily permutée & $\Delta$ \\
        \midrule
        0 & 0,6270 & 0,6265 & $+0{,}0004$ \\
        1 & 0,7597 & \textbf{0,7601} & $-0{,}0005$ \\
        3 & 0,8635 & \textbf{0,8644} & $-0{,}0009$ \\
        8 & 0,9374 & \textbf{0,9378} & $-0{,}0004$ \\
        \bottomrule
      \end{tabularx}
      \\[6pt]
      \statusbadge{failure}{l'alignement ne sert à rien}
    \end{column}
    \begin{column}{0.55\textwidth}
      \centering
      \compactpanel{\texttt{frangi\_sim}}{t2g_gain_riss_frangi.jpg}\hfill
      \compactpanel{\texttt{ofs}}{t2g_gain_riss_ofs.jpg}\hfill
      \compactpanel{\texttt{ofa}}{t2g_gain_riss_ofa.jpg}\hfill
      \compactpanel{\texttt{phase\_sym}}{t2g_gain_riss_phase.jpg}
      \\[4pt]
      {\scriptsize Image du \emph{plus grand gain} de la géométrie :
      \textbf{aucun canal ne désigne la fissure}.}
    \end{column}
  \end{columns}

  \vspace{0.26cm}
  \centering
  \statusbadge{inkblue}{l'adapter s'active pourtant : projections jusqu'à $1{,}79\times10^{-3}$}

  \source{Le déficit résiduel de \texttt{geo\_tol3} face à \texttt{tol3} vient
  des 290\,801 paramètres ajoutés, non d'un mauvais usage de la géométrie.}
\end{frame}

% -----------------------------------------------------------------------------
\begin{frame}{Et la galerie de gains trompe}
  \threewayrow{\texttt{CRACK500\_20160405\_171336}}{b2t_loss_crack171336}{t2g_gain_crack171336}{0,875 $\rightarrow$ 0,510 $\rightarrow$ 0,807}

  \vspace{0.22cm}
  \threewayrow{\texttt{Rissbilder\_for\_Florian}}{b2t_loss_riss}{t2g_gain_riss}{0,422 $\rightarrow$ 0,153 $\rightarrow$ 0,405}

  \vspace{0.14cm}
  \centering
  \statusbadge{failure}{la géométrie rend ce que la dilatation avait coûté, et reste sous la baseline}

  \vspace{0.10cm}
  {\fontsize{6.4}{7.4}\selectfont \colourkey}

  \source{Ces deux images sont les deux plus grands gains de la géométrie. Ce
  sont aussi les deux plus grandes pertes de \texttt{tol3}.}
\end{frame}

% -----------------------------------------------------------------------------
\begin{frame}{Pourquoi dilater}
  \centering
  \vspace{0.04cm}
  \synthpanel{parfait}{synth_perfect.png}\hfill
  \synthpanel{décalé 2 px}{synth_shifted.png}\hfill
  \synthpanel{2$\times$ trop large}{synth_wide.png}\hfill
  \synthpanel{rompu 4 px}{synth_break4.png}\hfill
  \synthpanel{rompu 20 px}{synth_break20.png}

  \vspace{0.20cm}
  \begin{columns}[T,onlytextwidth]
    \begin{column}{0.45\textwidth}
      \centering
      \small
      \setlength{\tabcolsep}{4pt}%
      \begin{tabularx}{\linewidth}{Xrr}
        \toprule
        IoU & strict & 3 px \\
        \midrule
        décalé de 2 px & 0,200 & \textbf{1,000} \\
        2$\times$ trop large & 0,429 & \textbf{1,000} \\
        rompu sur 4 px & 0,917 & \textbf{1,000} \\
        rompu sur 20 px & 0,583 & 0,708 \\
        \bottomrule
      \end{tabularx}
    \end{column}
    \begin{column}{0.51\textwidth}
      \small
      La vérité terrain \textbf{dilatée d'un pixel} n'obtient plus que
      \textbf{0,881} contre elle-même, et \textbf{0,799} à deux pixels.\\[4pt]
      Baseline, à 1 px de tolérance : la précision passe de 0,759 à
      \textbf{0,888}, le rappel seulement de 0,761 à 0,812.\\[4pt]
      \statusbadge{failure}{un décalage de 2 px est plus puni qu'une rupture de 20 px}
    \end{column}
  \end{columns}

  \source{La tolérance est une distance d'appariement : elle pardonne le
  placement et la largeur, et une rupture seulement si elle est plus courte que
  $2k$.}
\end{frame}

% -----------------------------------------------------------------------------
\begin{frame}{Ce que la dilatation gagne}
  \caserow{\texttt{cracktree200\_6774} --- réseau fin entièrement
    manqué}{b2t_gain_ct6774}{Baseline}{Avec
    \texttt{tol3}}{success}{0,000 $\rightarrow$ 0,219}

  \vspace{0.18cm}
  \caserow{\texttt{GAPS384\_0552} --- la contrepartie : débordement le long
    d'un bord clair}{b2t_loss_gaps552}{Baseline}{Avec
    \texttt{tol3}}{failure}{0,517 $\rightarrow$ 0,308}

  \vspace{0.20cm}
  \centering
  \small
  \setlength{\tabcolsep}{5pt}%
  \begin{tabular}{lrrrr}
    \toprule
    IoU tampon & strict & 1 px & 3 px & 8 px \\
    \midrule
    \ttfamily baseline & 0,6241 & 0,7407 & 0,8396 & 0,9213 \\
    \ttfamily tol3 & \textbf{0,6276} & \textbf{0,7586} & \textbf{0,8607} &
      \textbf{0,9346} \\
    \bottomrule
  \end{tabular}
\end{frame}

% -----------------------------------------------------------------------------
\begin{frame}{Suite}
  \small
  \begin{tabularx}{\linewidth}{c>{\raggedright\arraybackslash}X}
    \toprule
    \textbf{1} & \textbf{Adopter la perte tolérante à 3 px par défaut.} Gain
      acquis, aucun paramètre, aucun coût. \\
    \textbf{2} & \textbf{Arrêter le guidage géométrique sur Khánh Hà.} Le
      modèle est indifférent à l'alignement de l'évidence. \\
    3 & \textbf{Le tester où l'information existe} : FIND multimodal, portée et
      thermique, que SAM 2 ne peut pas voir. \\
    \bottomrule
  \end{tabularx}

  \vspace{0.30cm}
  \centering
  \statusbadge{failure}{à ne pas faire : plus de capacité, plus d'époques, un GNN}

  \vspace{0.20cm}
  \statusbadge{warning}{limites : 5 époques et non 20, pas d'IC95, \texttt{tol5} non exécuté}

  \source{Sur ce corpus --- monomodal visible, annotations épaisses, baseline
  supervisée sur le même domaine --- le Frangi généralisé n'apporte aucune
  information que la LoRA n'ait déjà apprise.}
\end{frame}

\end{document}
